% ── ABSTRACT, Foundations of Physics variant ──────────────────────────────
% CONFIRMED (Springer Nature's own submission-guidelines page, verified via
% web search): abstract must be 150-250 words.
% NOT independently confirmed: a "no equations, no citations in the
% abstract" rule. This draft avoids both anyway, as no-equations/
% no-citations abstracts are near-universal journal convention and a safe
% default -- but if you find Foundations of Physics explicitly permits
% inline citations/equations in abstracts, that constraint could be relaxed.
% Trimmed/rewritten from sections/00-abstract.tex (the Quantum-journal
% abstract, ~319 words with 2 display equations and 6 citations). Canonical
% main.tex abstract is untouched.

The Relational Coherence Framework (\RCF) does not alter the
predictions of quantum mechanics. Its contribution is explanatory and
operational: it recasts decoherence, property definiteness, and tunable
entanglement as aspects of a single coherence-budget process, in which
local relational potential is redistributed into environmental
correlation. This paper (i)~collects four standard identities of
pointer-basis coherence theory in the language of a normalized
Hilbert--Schmidt coherence quantifier grounded by einselection, adopting a
relational reading of quantum mechanics as the conceptual setting without
claiming it is singled out among the alternatives, and (ii)~derives an
allocation theorem for tunable Heisenberg-exchange gates that turns the
per-layer coherence cost into a compiler-relevant scheduling objective.

The central result, the Coherence-Budget Allocation Theorem, gives the
unique cost-minimizing exchange-angle schedule under Markovian
gate-time-proportional dephasing with a fixed entanglement target: the
schedule routes more entanglement onto quieter layers and skips layers
whose noise rate exceeds a global shadow price set by the budget.

A reference scheduler implementing this theorem beats both fixed-gate and
uniform-angle baselines in multi-layer density-matrix simulation, yielding
fidelity advantages of up to 0.040 against a concentrated full-entangler
baseline and 0.025 against a uniform-angle schedule on heterogeneous noise.
The framework's per-gate predictions remain consistent with experimental
continuous-tunable-gate results on superconducting hardware. The advantage
is contingent on direct exchange calibration; on platforms where the gate
must be decomposed into fixed native primitives, the advantage may shrink
or vanish.
